\subsection{Aperçu préliminaire : appréciation qualitative du pipeline}
\label{sub:exp0}

% ──────────────────────────────────────────────────────────────────────────────
\subsubsection{Intention et enjeux}

Avant d'engager un protocole d'évaluation quantitatif complet, il est
nécessaire de s'assurer que le pipeline d'extraction et d'encodage produit des
représentations cohérentes avec les propriétés théoriques attendues d'un
encodeur de type JEPA : 
\begin{enumerate}
  \item \emph{invariance à la compilation} — deux variants
compilés d'un même programme devraient être proches dans l'espace latent,
  \item \emph{séparabilité inter-programmes} — deux utilitaires de sémantique
distincte devraient être éloignés indépendamment du compilateur utilisé.
\end{enumerate}
Cette première expérimentation constitue une \emph{preuve de concept}
qualitative. 
Elle ne vise pas à mesurer une performance absolue, mais à valider
que la chaîne de traitement complète — compilation, extraction des
Bag-of-Paths, encodage par \texttt{elf64ctx}, comparaison et visualisation —
produit des signaux interprétables et alignés avec les hypothèses du modèle.
Son rôle est d'ancrer les expérimentations ultérieures : tout raffinage
méthodologique (affinement du corpus, modification de la stratégie de masquage,
introduction de métriques quantitatives) s'appuiera sur ce point de référence
qualitatif.

% ──────────────────────────────────────────────────────────────────────────────
\subsubsection{Méthodologie}

Le pipeline expérimental se décompose en quatre étapes automatisées,
implémentées dans le répertoire \texttt{experiments/0-firstOverview/}.

\paragraph{Étape 1 — Compilation des cibles.}
Quatre utilitaires GNU Coreutils (\texttt{cat}, \texttt{cp}, \texttt{ls},
\texttt{mv}) sont compilés selon une grille réduite de trois variantes :
\texttt{gcc~-O0}, \texttt{gcc~-O3} et \texttt{clang~-Os}. Chaque binaire est
strippé de ses symboles (\texttt{strip~--strip-all}), simulant ainsi les
conditions typiques d'analyse de code malveillant. Le résultat est un corpus
de $4 \times 3 = 12$ fichiers ELF64.

\paragraph{Étape 2 — Extraction des embeddings.}
L'outil \texttt{elf64ctx} est invoqué sur chaque binaire en mode encodage
(\texttt{-E}). Pour chaque fichier ELF, il extrait les \gls{cfg}
intra-procéduraux via \texttt{angr}, génère les
Bag-of-Paths par exploration DFS bornée (max.\ 500 chemins, max.\ 50 blocs
par chemin), projette les séquences VEX vers le vocabulaire canonique
$\mathcal{V}$, et encode chaque chemin via le Conv1DEncoder pré-entraîné. Les
embeddings, normalisés en norme $\ell_2$, sont sérialisés au format JSON
(sérialisation pickle + encodage Base64) pour chaque fichier ELF.

\paragraph{Étape 3 — Comparaison par paires.}
Pour chacun des $\binom{12}{2} = 66$ couples de binaires, \texttt{elf64ctx}
est invoqué en mode comparaison (\texttt{--compare}). La distance euclidienne
est calculée dans l'espace latent 256-dimensionnel entre toutes les paires
$(f_i,\,f_j)$ de fonctions issues des deux binaires, formant une matrice de
proximité de dimensions $|F_1| \times |F_2|$.

\paragraph{Étape 4 — Visualisation.}
L'option \texttt{--plot} génère une heatmap pour chaque matrice de distances,
en colormap \texttt{plasma\_r} : les teintes chaudes (proches de zéro)
indiquent une forte proximité sémantique, les teintes froides une forte
divergence. La Figure~\ref{fig:exp0_pipeline} résume l'enchaînement des étapes.

% ──────────────────────────────────────────────────────────────────────────────
% Figure : pipeline de l'expérimentation
% ──────────────────────────────────────────────────────────────────────────────
\begin{figure}[ht]
\centering
\begin{tikzpicture}[
  box/.style={draw, rounded corners=5pt, minimum width=2.6cm,
              minimum height=0.8cm, align=center, font=\small,
              line width=0.7pt},
  arr/.style={-Stealth, line width=0.8pt},
  lbl/.style={font=\tiny, gray!70!black},
]
  \node[box, fill=colCtx!15,  draw=colCtx!60]  (comp) at (0,0)
    {Compilation\\{\scriptsize\texttt{gcc}/\texttt{clang}}};
  \node[box, fill=colTgt!15,  draw=colTgt!60]  (emb)  at (3.4,0)
    {Extraction\\{\scriptsize\texttt{elf64ctx -E}}};
  \node[box, fill=colPred!15, draw=colPred!60] (cmp)  at (6.8,0)
    {Comparaison\\{\scriptsize\texttt{--compare}}};
  \node[box, fill=colLoss!15, draw=colLoss!60] (viz)  at (10.2,0)
    {Visualisation\\{\scriptsize\texttt{--plot}}};

  \draw[arr] (comp) -- (emb)  node[midway, above, lbl] {12 ELF};
  \draw[arr] (emb)  -- (cmp)  node[midway, above, lbl] {12 JSON};
  \draw[arr] (cmp)  -- (viz)  node[midway, above, lbl] {66 matrices};

  \node[lbl, below=6pt of comp] {4 utils $\times$ 3 variants};
  \node[lbl, below=6pt of emb]  {Bag-of-Paths + Conv1D};
  \node[lbl, below=6pt of cmp]  {distances $\ell_2$ dans $\mathbb{R}^{256}$};
  \node[lbl, below=6pt of viz]  {heatmaps \texttt{plasma\_r}};
\end{tikzpicture}
\caption{Pipeline de l'expérimentation 0. Les 12 binaires ELF produits à
  l'étape de compilation transitent successivement par l'extraction d'embeddings,
  la comparaison exhaustive par paires et la génération de heatmaps.}
\label{fig:exp0_pipeline}
\end{figure}

% ──────────────────────────────────────────────────────────────────────────────
\subsubsection{Résultats — Analyse des heatmaps}

Les figures ci-après présentent cinq heatmaps représentatives sélectionnées
parmi les 66 générées. Elles couvrent les deux régimes d'intérêt : les
\emph{paires positives} (même programme, compilations distinctes) et les
\emph{paires négatives} (programmes différents), ainsi qu'un cas mixte
combinant les deux sources de divergence.

\paragraph{Paires positives — résilience à la variation de compilation.}
Les Figures~\ref{fig:exp0_cat_gcc0} et~\ref{fig:exp0_cat_gcc3} mettent en
regard \texttt{cat} compilé avec \texttt{clang~-Os} et, respectivement,
\texttt{gcc~-O0} et \texttt{gcc~-O3} : même utilitaire, compilateurs et
niveaux d'optimisation distincts. Dans les deux cas, l'échelle des distances
se maintient dans l'intervalle $[0,\,1.0]$ et les matrices révèlent des
\textbf{zones localement chaudes} ($\approx 0.2$--$0.4$), correspondant à des
sous-ensembles de fonctions que le modèle a appris à rapprocher malgré la
transformation compilateur. Ces îlots de chaleur attestent d'une
\emph{invariance sémantique partielle} acquise lors du pré-entraînement : les
Bag-of-Paths preservent suffisamment de structure comportementale en
\gls{vexir} pour que des fonctions compilées différemment partagent un
voisinage dans l'espace latent.

On notera l'absence d'une diagonale chaude explicite, ce qui est attendu :
les adresses de fonctions ne sont pas ordonnées de façon cohérente entre deux
binaires distincts, et certaines fonctions peuvent être inlinées, fusionnées
ou scindées selon le niveau d'optimisation.

\paragraph{Paires négatives — séparabilité inter-utilitaires.}
Les Figures~\ref{fig:exp0_cat_cp_cl} et~\ref{fig:exp0_cat_mv_cl} opposent
\texttt{cat~(clang~-Os)} respectivement à \texttt{cp~(clang~-Os)} et
\texttt{mv~(clang~-Os)} : utilitaires différents, compilateur identique. Ces
paires affichent des matrices uniformément froides dans l'intervalle $[0.6,\,
1.0]$, illustrant que le modèle positionne des programmes de sémantique
distincte dans des régions éloignées de l'espace latent, même lorsque le
compilateur est identique.

\paragraph{Cas mixte — divergence super-additive.}
La Figure~\ref{fig:exp0_cat_cp_gcc0} compare \texttt{cat~(clang~-Os)} à
\texttt{cp~(gcc~-O0)} : deux utilitaires \emph{différents} compilés avec des
\emph{compilateurs et optimisations différents}. L'échelle des distances atteint
$[0,\,3.0]$, soit un écart trois fois supérieur à celui observé dans les autres
comparaisons. Ce résultat suggère que la combinaison de deux sources de
divergence — sémantique fonctionnelle et style compilateur — amplifie la
séparation dans l'espace latent de façon \textbf{super-additive}, propriété
désirable pour une utilisation en détection de similarité binaire.

% ──────────────────────────────────────────────────────────────────────────────
% Heatmaps : 2 paires positives
% ──────────────────────────────────────────────────────────────────────────────
\begin{figure}[ht]
  \centering
  \begin{minipage}[t]{0.47\textwidth}
    \centering
    \includegraphics[width=\linewidth]{%
      ../../experiments/0-firstOverview/data/comparisons/cat_clang_Os_vs_cat_gcc_O0.png}
    \caption{\textbf{Paire positive.} \texttt{cat~(clang~-Os)} vs
      \texttt{cat~(gcc~-O0)}. Même sémantique, compilateurs distincts.
      Distances dans $[0,\,1.0]$ ; îlots chauds ($\approx\!0.2$--$0.4$)
      témoignant d'une invariance partielle à la compilation.}
    \label{fig:exp0_cat_gcc0}
  \end{minipage}
  \hfill
  \begin{minipage}[t]{0.47\textwidth}
    \centering
    \includegraphics[width=\linewidth]{%
      ../../experiments/0-firstOverview/data/comparisons/cat_clang_Os_vs_cat_gcc_O3.png}
    \caption{\textbf{Paire positive.} \texttt{cat~(clang~-Os)} vs
      \texttt{cat~(gcc~-O3)}. Niveau d'optimisation agressif côté GCC.
      Motif analogue : l'invariance persiste malgré l'inlining intensif de
      \texttt{-O3}.}
    \label{fig:exp0_cat_gcc3}
  \end{minipage}
\end{figure}

% ──────────────────────────────────────────────────────────────────────────────
% Heatmaps : 2 paires négatives (même compilateur)
% ──────────────────────────────────────────────────────────────────────────────
\begin{figure}[ht]
  \centering
  \begin{minipage}[t]{0.47\textwidth}
    \centering
    \includegraphics[width=\linewidth]{%
      ../../experiments/0-firstOverview/data/comparisons/cat_clang_Os_vs_cp_clang_Os.png}
    \caption{\textbf{Paire négative.} \texttt{cat~(clang~-Os)} vs
      \texttt{cp~(clang~-Os)}. Compilateur identique, sémantique distincte.
      Matrice uniformément froide dans $[0.6,\,1.0]$ : bonne séparabilité
      inter-utilitaires.}
    \label{fig:exp0_cat_cp_cl}
  \end{minipage}
  \hfill
  \begin{minipage}[t]{0.47\textwidth}
    \centering
    \includegraphics[width=\linewidth]{%
      ../../experiments/0-firstOverview/data/comparisons/cat_clang_Os_vs_mv_clang_Os.png}
    \caption{\textbf{Paire négative.} \texttt{cat~(clang~-Os)} vs
      \texttt{mv~(clang~-Os)}. Même compilateur, programmes distincts.
      Résultat cohérent avec la Figure~\ref{fig:exp0_cat_cp_cl} : le modèle
      maintient la séparation sur plusieurs utilitaires négatifs.}
    \label{fig:exp0_cat_mv_cl}
  \end{minipage}
\end{figure}

% ──────────────────────────────────────────────────────────────────────────────
% Heatmap : cas mixte super-additif
% ──────────────────────────────────────────────────────────────────────────────
\begin{figure}[ht]
  \centering
  \includegraphics[width=0.55\linewidth]{%
    ../../experiments/0-firstOverview/data/comparisons/cat_clang_Os_vs_cp_gcc_O0.png}
  \caption{\textbf{Cas mixte — divergence super-additive.}
    \texttt{cat~(clang~-Os)} vs \texttt{cp~(gcc~-O0)}. Sémantique différente
    \emph{et} compilateurs différents. L'échelle des distances atteint
    $[0,\,3.0]$, soit un facteur $\approx\!3\times$ supérieur aux paires
    négatives à compilateur identique. La combinaison des deux sources de
    divergence amplifie la séparation dans l'espace latent de façon
    super-additive.}
  \label{fig:exp0_cat_cp_gcc0}
\end{figure}

% ──────────────────────────────────────────────────────────────────────────────
\paragraph{Synthèse.}
Ces observations valident qualitativement deux propriétés attendues d'un
encodeur JEPA entraîné sur du VEX IR : (i) une \textbf{invariance partielle à
la compilation} pour un même programme, se manifestant par le rapprochement
des représentations dans les paires positives ; (ii) une \textbf{séparabilité
inter-programmes}, renforcée lorsque les sources de variabilité s'accumulent.
Les limites — absence de diagonale franche, proximité résiduelle sur certaines
paires négatives à compilateur identique — sont cohérentes avec un modèle
entraîné en preuve de concept sur un corpus de taille restreinte.

% ──────────────────────────────────────────────────────────────────────────────
\subsubsection{Limites et perspectives}

Cette expérimentation appelle plusieurs observations critiques qui motivent les
expérimentations futures.

\paragraph{Absence de référence assembleur.}
La principale faiblesse de cet essai est l'absence de \textbf{mise en
correspondance directe entre les représentations assembleur et les embeddings
produits}. Les heatmaps permettent d'observer la structure globale de l'espace
latent, mais elles ne permettent pas de répondre à la question suivante : les
fonctions que le modèle juge proches sont-elles effectivement similaires au
niveau du code désassemblé ? À l'inverse, les régions froides correspondent-elles
à des fonctions dont l'assembleur diffère structurellement ?

En l'absence d'un \textbf{ancrage corrélationnel visuel} — par exemple une
visualisation côte à côte de l'alignement fonctionnel assembleur et de la
heatmap de distances — il est impossible de distinguer une invariance
sémantique apprise d'un artefact d'encodage (effondrement de représentations
ou biais de distance systématique). Cet ancrage constitue le prérequis
méthodologique des expérimentations futures : avant de comparer des métriques
de similarité agrégées, il est nécessaire de s'assurer que les voisinages
latents correspondent à des voisinages sémantiques vérifiables dans
l'assembleur.

\paragraph{Absence de métriques quantitatives.}
Les heatmaps sont interprétées qualitativement. Aucune métrique formelle —
précision de récupération (\textit{recall@k}), AUC sur un classifieur
paires-positives/négatives, \textit{mean average precision} — n'est calculée.
L'évaluation reste donc subjective et ne permet pas de comparaison rigoureuse
entre configurations d'entraînement.

\paragraph{Corpus minimal.}
Seuls quatre utilitaires et trois variantes de compilation sont utilisés
($12$ binaires, $66$ paires). Ce sous-corpus ne couvre pas la diversité
sémantique du dataset complet ($1\,060$ ELF) et peut exhiber des biais de
sélection. En particulier, \texttt{cat}, \texttt{cp}, \texttt{ls} et
\texttt{mv} partagent des primitives de bas niveau (appels à
\texttt{<API\_WRITE>}, \texttt{<SYSCALL\_READ>}) qui pourraient introduire une
proximité résiduelle non représentative.

\paragraph{Perspectives.}
Ces limites définissent le programme des expérimentations suivantes. La
priorité est l'introduction d'un protocole d'évaluation quantitatif (métriques
de récupération) couplé à un ancrage assembleur : pour un ensemble de paires
de fonctions, comparer leur distance dans l'espace latent à une distance de
référence calculée sur leur représentation \gls{vexir} tokenisée (similarité
de Jaccard sur les n-grammes de tokens, par exemple). Cet ancrage permettra de
mesurer dans quelle mesure l'espace latent capture des distances sémantiquement
significatives, et non de simples corrélations de surface liées au compilateur
ou à la taille des fonctions.
